% !TEX root = ../main.tex
% Student declaration: the statement of Unisa's Procedures for master's and doctoral degrees
% (Senate, 28 October 2026), 7.19, with the title of the thesis. The candidate signs and dates
% it before submission, after the originality report is in hand.

\clearpage
\phantomsection
\chapter*{선언}
\label{ch:declaration}

% Set a little closer than the text so that the statement and the signature stand on one page.
\begingroup\setstretch{1.35}
\noindent 이름: 권태홍 (Taehong Kwon) \hfill 학번: 28576810\\
학위: 컴퓨터과학 박사

\smallskip
\noindent ERC-20 허락 화살표를 페이지랭크에 넣어 온체인 지갑 평판 점수를 더 좋게 만들기

\smallskip
저는 위의 논문이 제가 직접 한 일이라고 선언해요. 그리고 제가 쓰거나 그대로 따온 자료는 모두 빠짐없는 참고문헌으로 표시하고 출처를 밝혔다고 선언해요. 또 인공지능(AI)을 쓴 일은 어떤 것이든 빠짐없이 밝혔다고 선언해요.

또 저는 이 논문을 Unisa(남아프리카대학교)가 인정하는 알맞은 독창성 검사 시스템(다른 글을 베끼지 않았는지 살펴보는 시스템)에 넣어 검사받았다고 선언해요. 그리고 그 결과가 대학교가 받아들이는 독창성 기준 안에 든다고 선언해요.

또 저는 이 연구를 전에 심사받으려고 다음 두 곳에 낸 적이 없다고 선언해요. 연구 전체도, 그 일부도 마찬가지예요. 다른 자격을 받으려고 Unisa에 낸 적이 없고, 다른 고등 교육 기관(대학 같은 곳)에 낸 적도 없어요.

또 저는 이 논문을 준비하면서 인공지능(AI) 도구를 쓴 곳에 대해 다음을 선언해요. 생성형 도구든 아니든 마찬가지예요. AI 도구는 윤리에 맞고 허용되는 도움을 받는 데에만 썼어요. 그 도움은 글을 더 분명하게 하고 일을 더 효율적으로 하려는 것이었고, 그러는 동안에도 학문적 진실성과 지은이의 책임을 지켰어요. 그렇게 쓴 일은 빠짐없이 밝혔어요. 그리고 AI 도구를 쓴 일은 제가 직접 한 독창적인 연구, 제가 스스로 한 비판적인 생각과 분석, 지은이로서 지는 책임을 대신하지 않아요.

저는 AI를 쓴 일을 밝히지 않는 것, 표절(남의 글을 자기 글인 것처럼 가져다 쓰는 일), 학문적 진실성(연구와 글쓰기를 정직하게 하는 태도)이 모자란 것은, 하나만 있어도 여럿이 함께 있어도, Unisa의 정책에 따라 학문적 부정행위가 될 수 있다는 것을 알고 있어요.

\smallskip
\noindent\textbf{생성형 AI 도구를 쓴 일.} 생성형 AI 도구(글이나 코드를 만들어 주는 인공지능 도구)인 Claude(Anthropic)가 다음 일을 도왔어요. 문장 다듬기, 글의 짜임을 고치고 내용이 앞뒤로 맞는지 서로 맞춰 보기, 분석 스크립트(분석에 쓰는 짧은 프로그램)의 일부를 쓰고 시험해 보기, 그리고 참고할 글을 찾고 확인하기예요. 연구 질문, 연구 설계, 분석, 결론은 모두 제 것이에요.

\vspace{1.8em}
\noindent\rule{6cm}{0.4pt}\hfill\rule{4cm}{0.4pt}\\
서명 \hfill 날짜
\par\endgroup
